\section{Data for numerical and graphical Output} \label{sec2.11}

\medskip

\textbf{General remarks}

\medskip %

The input flags in this block control all the numerical and graphical
output of an EIRENE run. This comprises diagnostics during the
initialization phase (e.g. 2d and 3d geometry plots) as well as selected
test particle histories printed and plotted during their generation.
All numerical output of a run is arranged in so called ``tallies".
There are volume averaged tallies, surface averaged tallies
(``surface crossing tallies") and global tallies. These latter
tallies are derived from the former ones by integration over the total
computational volume, or over all non-transparent
surfaces, respectively.

Each volume averaged tally is an array TALLYV(I) of NSBOX data,
comprising the responses for the respective detector function $g_V
$ in each zone of the computational mesh (see
\cref{sec2.2}, ``General Remarks").

Each surface averaged tally is an array TALLYS(I) of NLIMPS data,
comprising the responses for the respective detector function
$g_S $ integrated over
each (non default standard or additional) surface.

Each volume or surface tally is characterized by its
index ITALV or ITALS respectively, augmented by a first (species) index
in some cases. The first index, if there is one, is always referred to
as ``species index" (by abuse of language), even for additional tallies
ADDV or ALGV in which cases it might have a different meaning.

As for the volume tallies there are NTALV preprogrammed detector
functions, up to NADV user supplied detector functions for
track-length estimators (tally no. ITALV) and NCLV user supplied
detector functions for collision estimators (tally no. ITALV).
Furthermore there are NSNV snapshot tallies (tally no. ITALV),
NCPV tallies for coupling with plasma background codes (tally no.
ITALV) and NBGV tallies for the (non-linear) BGK collision
integrals, needed for the iteration procedure (tally no. ITALV).
Up to NALV algebraic expressions in these tallies (tally no.
ITALV) can be stored.

There is also one tally no. 0, which is evaluated in subroutine
TALUSR (\cref{sec3.7}). This tally is printed immediately
after evaluation, i.e., it is not put onto any storage (unless
explicitly done so in subroutine TALUSR).

The input volume tallies for the background medium are listed in
\cref{tab1}. They are selected for printout or graphical
output by the \textbf{negative} value of their tally number. E.g.,
the choice ITALV = 5 selects tally no. 5 from \cref{tab2},
and the choice ITALV = -5 selects tally no. 5 from
\cref{tab1}.

Furthermore, there are NTALS  preprogrammed detector functions for
surface averaged tallies, and there may be up to NADS additional
user supplied surface averaged tallies (tally no. ITALS).

The preprogrammed tallies and their default units are given in
\cref{tab2} (volume averaged tallies) and in \cref{tab3}
(surface averaged tallies).

\cref{tab1} lists the NTALI  (volume-) tallies which are
computed from the input data in block 5 and which describe the
background medium (plasma). By an abuse of language the cell
volumes ($cm^{-3}$) stored in the array VOL(I) are also referred
to as an input volumetric tally (no. ITALI=13 in
\cref{tab1}).

\medskip

\textbf{The Input Block}

\medskip

\textsl{ * 11 Data for numerical and graphical output}

\textsl{* 11A Block for numerical output written on unit 6 (IUNOUT)}
{\small
\begin{verbatim}
 TRCPLT  TRCHST  TRCNAL  TRCREA  TRCSIG
 TRCGRD  TRCSUR  TRCREF  TRCFLE  TRCAMD
 TRCINT  TRCLST  TRCSOU  TRCREC  TRCTIM
 TRCBLA  TRCBLM  TRCBLI  TRCBLP  TRCBLE
 TRCBLPH TRCTAL  TRCOCT  TRCCEN  TRCRNF
 TRCHKTIM
 TRCSRC(J), J= 0, NSTRAI
 NVOLPR, NSPCPR
 DO 101 J=1,NVOLPR
         NPRTLV(J) NFLAGV(J) NSPEZV(J,1) NSPEZV(J,2) NTLVFL(J)
 101 CONTINUE
 NSURPR
 DO 102 J=1,NSURPR
    NSRF NPRTLS(J) NFLAGS(J) NSPEZS(J,1) NSPEZS(J,2) NTLSFL(J)
 102 CONTINUE
\end{verbatim}
}
optional: unless block 11B starts:
\begin{verbatim}
NLTVOUT
  NUMTAL(J),J=1,NLTVOUT  (12 values per card, FORMAT 12I6)
NLTSOUT
  NUMTAL(J),J=1,NLTSOUT  (12 values per card, FORMAT 12I6)
\end{verbatim}

 \textsl{ * 11B Data for graphical output}

 \textsl{ * 11B.1 Data for 2D and 3D Geometry plot}
{\small
\begin{verbatim}
 PL1ST     PL2ND     PL3RD     PLADD  PLHST
 PLCUT(1)  PLCUT(2)  PLCUT(3)  PLBOX  PLSTOR
 PLNUMV    PLNUMS    PLARR
 NPLINR NPLOTR NPLDLR NPLINP NPLOTP NPLDLP NPLINT NPLOTT NPLDLT
 DO 1140 J=1,5
    PL3A(J) TEXTLA(J) IPLTA(J) (IPLAA(J,I) IPLEA(J,I) I=1,IPLTA(J))
 1140 CONTINUE
 DO 1141 J=1,3
    PL3S(J) TEXTLS(J) IPLTS(J) (IPLAS(J,I) IPLES(J,I) I=1,IPLTS(J))
 1141 CONTINUE
 CH2MX  CH2MY         CH2X0  CH2Y0  CH2Z0
 CH3MX  CH3MY  CH3MZ  CH3X0  CH3Y0  CH3Z0
 ANGLE1  ANGLE2  ANGLE3
 I1TRC  I2TRC  (ISYPLT(J)  J=1,8)  ILINIE
\end{verbatim}
}
 \textsl{ * 11B.2 Data for plots of volume tallies}
{\small
\begin{verbatim}
 NVOLPL
 PLTSRC(J), J= 0, NSTRAI
 DO 116 N=1,NVOLPL
 optional: one text card starting with *...... to label the  plot
    NSPTAL(N)
    PLTL2D(N) PLTL3D(N) PLTLLG(N) PLTLER(N)
    TALZMI(N) TALZMA(N) TALXMI(N) TALXMA(N) TALYMI(N) TALYMA(N)
  IF (PLTL2D(NTL)) THEN
    LHIST2(N) LSMOT2(N)
    DO 118 I=1,NSPTAL(N)
       ISPTAL(N,I) NPTALI(N,I) NPLIN2(N,I) NPLOT2(N,I)
       NPLDL2(N,I)
    118 CONTINUE
  ENDIF
  IF (PLTL3D(N)) THEN
    LHIST3 LCNTR3 LSMOT3 LRAPS3 LVECT3 LRPVC3 LRAPS3D LR3DCON
    LPRAD LPPOL LPTOR
    DO 119 I=1,NSPTAL(N)
       ISPTAL(N,I)  NPTALI(N,I) IPROJ3(N,J)
       NPLI13(N,I)  NPLO13(N,I)
       NPLI23(N,I)  NPLO23(N,I) IPLANE(N,I)
    119 CONTINUE
    TALW1(N)  TALW2(N)  FCABS1(N)  FCABS2(N) RAPSDEL(N)
  ENDIF
 116 CONTINUE
\end{verbatim}
}

\medskip

\textbf{Meaning of the Input Variables for numerical and graphical Output}

\medskip

\textsl{* 11A. Block for numerical output written on unit 6}
\begin{list}{ }{}
\item[\textbf{TRCPLT }]
Trace-back from plot routines
\item[\textbf{TRCHST }]
Printout of trajectories of selected test particle histories into
geometry plots.
These histories are selected by the flags I1TRC and I2TRC, see below,
sub-block 11B.2.
\item[\textbf{TRCNAL }]
Trace-back for non-analog methods, i.e. splitting surfaces,
suppression of absorption, weighted post-collision species
sampling etc.
\item[\textbf{TRCMOD }]
Trace-back from routines for iterative mode (MOD\_TMSTEP, MODBGK, and problem
specific routines called from MODUSR).
\item[\textbf{TRCSIG }]
Trace-back from post processing line integral diagnostics block
DIAGNO, line of sight integration routines.
\item[\textbf{TRCGRD }]
Printout of ``standard mesh surface data"
\item[\textbf{TRCSUR }]
Printout of data for ``additional surfaces"
\item[\textbf{TRCREF }]
Printout of reflection model related data.
In particular a list of all non-perfect recycling surfaces is printed,
i.e., a list of surfaces for which there is some absorption at least for
one incident particle species.
\item[\textbf{TRCFLE }]
Trace-back from subroutines WRSTRT, WRGEOM, WRPLAS (writing on and
reading from the dump files FT10, FT11, FT12, FT13 etc.
\item[\textbf{TRCAMD }]
Trace-back from atomic and molecular data routines

XSECTA, XSECTM, XSECTI
\item[\textbf{TRCINT }]
Traceback from user specified interfacing routine INFCOP,
e.g. of data related to coupling of EIRENE to other codes.
\item[\textbf{TRCLST }]
Printout of information during the last history of each stratum.
E.g., sampling efficiencies, and other accumulated information,
which is only available in the history generation routines during
particle tracing.
\item[\textbf{TRCSOU }]
Traceback from primary source sampling routines, e.g. from LOCATE,
SAMPNT, SAMLNE, SAMSRF, SAMVOL and SAMUSR
\item[\textbf{TRCREC }]
Printout of
EIRENE recommendations for next run on same case:

A) stratified source sampling, at present:
for proportional allocation of weights.

B) weight windows, at present: to be written
\item[\textbf{TRCTIM }]
Printout cpu-time information for each history
\item[\textbf{TRCBLA}]
Global particle and energy balance for atoms is printed.
\item[\textbf{TRCBLM }]
Global particle and energy balance for
molecules is printed.
\item[\textbf{TRCBLI }]
Global particle and
energy balance for test ions is printed.
\item[\textbf{TRCBLP }]
Global particle and energy balance for bulk ions is printed.
\item[\textbf{TRCBLE }]
Global particle and energy balance for
electrons is printed.
\item[\textbf{TRCBLPH }]
Global particle and
energy balance for photons is printed.
\item[\textbf{TRCTAL }]
Print list of activated and de-activated tallies. The default settings eliminate some tallies from storage and
estimators which are likely to be irrelevant in a particular run, e.g. all photon related tallies are deactivated automatically
if no radiation transfer calculation is included in a run.
These default settings are overruled by the flags NTLVOUT and NTLSOUT at the end of this sub-block 11A, see below.
\item[\textbf{TRCOCT }]
to be written: printout from octree geometry optimization procedure
\item[\textbf{TRCCEN }]
printout from census array stored in time dependent mode: species and stratum resolved census fluxes
\item[\textbf{TRCRNF }]
printout from random number generation routines: setting seed, etc., mostly for testing ``correlated sampling'' options (flag: NLCRR).
\item[\textbf{TRCHKTIM }]
printout of cpu usage during particle tracing and scoring (excluding overhead), separated by particle type, species and strata.
\item[\textbf{TRCSRC(ISTRA),ISTRA=0,NSTRAI }] ~

The selected global, volumetric and
surface crossing tallies are printed for those
strata for which TRCSTR(ISTRA) = .TRUE.. In case
TRCSTR(0) = .TRUE., the results after
summation over all strata is printed for
these tallies.

Note: printout for individual strata is only possible if the
data for strata have been saved on file (NFILE-N=1,2 option, input block 1).
Otherwise only the last stratum currently on storage arrays (mostly: sum over strata) is available.
\item[\textbf{NVOLPR }]
Total number of volume averaged tallies to be printed.
\item[\textbf{NSPCPR }]
Flag for printout of surface or cell based spectra (defined in input block 10F)
If NSPCPR $> 0$: spectra are printed on output stream fort.(20+ioff), name: spectra.out
\item[\textbf{NPRTLV }]
Index of the tally to be printed (first column in \cref{tab1} or
\cref{tab2}). If the tally has a species index, it is printed for all
species,
for which the integral of the tally over the computational area is
nonzero. Otherwise the statement

ZERO INTEGRAL FOR SPECIES ISPZ = ......

is printed below the header for this tally.
The term ``species index" is used here in a more general sense  for the first index (if any)
for any given tally. In case of some ``additional tallies" or other tallies,
in which the first index does not label a particle species but something
else, this term ``species index" then is to be understood in a more
general sense.

If NPRTLV = 0, then the user defined post processing routine TALUSR
is called (\cref{sec3.7}).
\item[\textbf{NFLAGV }]
Flag to specify the level of printout.
\begin{list}{ }{}
\item[{$=-1$ }]
print only header
\item[{$=0$ }]
additionally: print global quantities (total and block averages)
\item[{$=1$ }]
additionally: print 1D profiles (if any), averaged over all (if any) higher dimensions
\item[{$=2$ }]
additionally: print 2D profiles (if any), averaged over all (if any) higher dimensions
\item[{$=3$ }]
additionally: print 3D profiles (if any), averaged over all (if any) higher dimensions
\item[{$=4$ }]
print 3D profiles, but no lower dimensional averages
\end{list}
\item[\textbf{NSPEZV }]
If NSPEZV(..,1) $\ne 0$, then this tally is printed only for the
species
index range:

$I1 = NSPEZV(...,1)$ to $I2 = MAX(I1,NSPEZV(...,2))$,

rather than for all species relevant for the specified tally.
\item[\textbf{NTLVFL }]
In addition to the standard output stream fort.IUNOUT, this tally is
also printed onto output stream fort.NTLVFL, for all species selected,
and also the corresponding standard deviations are printed, if
available.
The format for this extra output stream is specified in subroutine
PRTTAL, in code segment EIRASS.
\end{list}
The next paragraph describes printout for surface averaged tallies. Distinct from the
volume averaged tallies, where selected tallies (NPRTLV) are printed for all the computational domain,
no particular surface tallies are selected, but always all kinds of surface fluxes (tallies)
are printed (particle fluxes, incident, emitted, energy fluxes,
sputtered fluxes,....).
A choice is made now instead with respect to the spatial region:
the particular surfaces (internal EIRENE surface
numbers) are selected here by input flag NSRF.

\begin{list}{ }{}
\item[\textbf{NSURPR }]
Total number of surfaces, for which surface averaged tallies are to be printed.
\item[\textbf{NSRF }]
Index of the surface
to be printed. By default all tallies listed in \cref{tab3,tab3a,tab3old} (depending on code version)
are printed for this surface.
\item[\textbf{NPRTLS(J) NFLAGS(J) NSPEZS(J,1) NSPEZS(J,2)}]
These next flags are only needed for those surfaces, for which a
further spatial resolution within one surface is provided (this is
enabled by providing storage through setting the flag NGSTAL = 1 in
input block 1). Their meaning then corresponds to the meaning of
flags NPRTLV(J), NFLAGV(J), NSPEZV(J,1), NSPEZV(J,2) for volume
tallies, respectively. If the NPRTLS,... flags are not specified
(i.e.: default=0), then no spatially resolved surface tallies are
printed.
\item[\textbf{NTLSFL(J)}] The spatially resolved tallies (if any) and/or the flux
energy spectra (if any, see input block 10F) are printed on
additional output stream fort.NTLSFL(J) for this surface.
\end{list}

The next, optional input cards can be used to overrule the default choices of tallies which are available in a run.
The logical flag TRCTAL described above can be used to print a list of all activated and deactivated tallies.
\begin{list}{ }{}
\item[\textbf{NTLVOUT }]
total number of volume averaged tallies to be explicitly abandoned or enabled
\item[\textbf{NUMTAL(J) }]
Number of a tally from \cref{tab2}, e.g. NUMTAL(J)=1 for neutral atom density tally PDENA.
The tally with this number NUMTAL(J) is explicitly enabled.
With an additional negative sign this tally is removed from the run (and the balances), e.g. NUMTAL(J)=-2
would remove the evaluation (and storage) of tally PDENM from this run.
\item[\textbf{NTLSOUT }]
to be written
\end{list}


\textsl{ * 11B.1 Data for 2D plot of geometry}

\begin{list}{ }{}
\item[\textbf{PL1ST }]
plot the x- (radial) standard grid surfaces into a 2D geometry plot
\item[\textbf{PL2ND }]
plot the y- (poloidal) standard grid surfaces into a 2D
geometry plot
\item[\textbf{PL3RD }]
plot the z- (toroidal) standard grid surfaces into a 2D
geometry plot
\item[\textbf{PLADD }]
plot the additional surfaces into a 2D geometry plot
\item[\textbf{PLHST }]
Plot the track of some selected test particle histories into the
geometry plot (2D or 3D).
\item[\textbf{PLCUT }]
Flag for the choice of a plane, in which the 3 dimensional
geometrical configuration is plotted in case of a 2D geometry plot.
This plane may be defined by either x = CONST, y = CONST or by
z = CONST.
\begin{list}{ }{}
\item[{PLCUT(1) }] = .TRUE.

x = CONST ; plotting plane is the yz-plane.

y is the ordinate, z is the abscissa
\item[{PLCUT(2) }] = .TRUE.

y = CONST ; plotting plane is the xz plane.

z is the ordinate, x is the abscissa
\item[{PLCUT(3) }] = .TRUE.

z = CONST ; plotting plane is the xy plane.

y is the ordinate, x is the abscissa
\end{list}
\item[\textbf{PLBOX }] plot box defined by surface inequalities
in addition to valid part of surface
\item[\textbf{PLSTOR }]
produce file of coordinates along 2D projections of ``additional
surfaces" for later use in some graphics (``PATRAN format", also
for RAPS-graphics, see below, sub-block: 11B3). These coordinates
are stored on arrays XPL2D, YPL2D in subroutine STCOOR called from
subroutine PLTADD.
\item[\textbf{PLNUMV }] print cell number
into standard mesh cells
\item[\textbf{PLNUMS }] print numbers
near additional surfaces
\item[\textbf{PLARR }] plot arrows to
indicate the outer surface normal vector (only in LEVGEO=3, partly in LEVGEO=2, for non-default standard surfaces (input block 3a).
Not functional (any more?) for additional surfaces.
\item[\textbf{PLIDL }] print output files for further processing
in external  tools (e.g. for FZJ proprietary IDL graphics tool).
The grid information as well as all input and output tallies (excluding the de-activated ones)
are printed on files, one file per tally. See calls to routines OUTIDL.... from main program EIRENE.
\item[\textbf{CH2MX }] horizontal
half width of 2D plot window (cm)
\item[\textbf{CH2MY }]
vertical half width of 2D plot window (cm)
\item[\textbf{CH2X0
}] horizontal co-ordinate of midpoint of 2D plot window (cm).
\item[\textbf{CH2Y0 }] vertical co-ordinate of midpoint of 2D
plot window (cm).
\item[\textbf{CH2Z0 }] distance CONST of
plotting plane to origin.
\item[\textbf{NPLINR NPLOTR NPLDLR }]
~

radial standard surfaces with labels

IR1ST = NPLINR, NPLOTR, NPLDLR

are plotted.
\item[\textbf{NPLINP NPLOTP NPLDLP }] ~

poloidal standard surfaces with labels

IP2ND= NPLINP, NPLOTP, NPLDLP

are plotted.
\item[\textbf{NPLINT NPLOTT NPLDLT }] ~

toroidal standard surfaces with labels

IT3RD= NPLINT, NPLOTT, NPLDLT

are plotted.
\end{list}

\textsl{ * 11B.2 Data for 3D plot of geometry}

\begin{list}{ }{}
\item[\textbf{CH3MX }]
half width of plot chamber in x- direction, used for
3D geometry plot
\item[\textbf{CH3MY }]
half width of plot chamber in y- direction, used for
3D geometry plot
\item[\textbf{CH3MZ }]
half width of plot chamber in z- direction, used for
3D geometry plot
\item[\textbf{CH3X0 }]
x-co-ordinate of midpoint of plot-chamber in user co-ordinates, 3D
geometry plot only
\item[\textbf{CH3Y0 }]
y-co-ordinate of midpoint of plot-chamber in user co-ordinates, 3D
geometry plot only
\item[\textbf{CH3Z0 }]
z-co-ordinate of midpoint of plot-chamber in user co-ordinates, 3D
geometry plot only
\item[\textbf{ANGLE1, ANGLE2 }] ~

First and second viewing angle for 3D geometry plot
\end{list}

The following 5 cards specify up to 5 groups of additional surfaces
to be plotted on the 3D geometry plot.
Each group may consist of subgroups, and each subgroup is defined
by an interval of additional surface labeling indices ILIMI,
from the full range of all (i.e NLIMI) additional surfaces.

\begin{list}{ }{}
\item[\textbf{PL3A }]
logical flag, indicating if the additional surfaces specified
in this card are to be plotted or not. If PL3A=.FALSE., the rest
of this card is irrelevant
\item[\textbf{TEXTLA }]
text written onto plot, characterizing this group of additional
surfaces
\item[\textbf{IPLTA }]
number of different subgroups of additional surfaces comprising
this group
\item[\textbf{IPLAA, IPLEA }] ~

each \hfill subgroup \hfill consists \hfill of \hfill additional \hfill
surfaces \hfill ranging \hfill from \hfill no.

IPLAA(J) to IPLEA(J), J=1,IPLTA
\end{list}

The following 3 cards reading PL3S,.....
specify a group of standard mesh surfaces
to be plotted on the 3D geometry plot.
The meaning of the flags in each card is the same as above for the
additional surfaces.
The first card refers to the first (x or radial) standard mesh,
the second card refers to the second (y or poloidal) standard mesh,
the third card refers to the third (z or toroidal) standard mesh.
Each group may consist of subgroups, and each subgroup is defined
by an interval of standard surface labeling indices ISURF,
from the full range of all (i.e NR1ST, NP2ND or NT3RD, resp.)
standard surfaces.

\begin{list}{ }{}
\item[\textbf{I1TRC }]
Number of first particle history
to be traced in printout and/or 2D or 3D plot
\item[\textbf{I2TRC }]
Number of last particle history
to be traced in printout and/or 2D or 3D plot
\item[\textbf{ISYPLT(J), J=1,8 }] ~

Indices to plot symbols along the particle tracks for different
events. Up to 8 different events can be picked in any order in the
array ISYPLT.
\begin{list}{ }{}
\item[{0 }] no symbol
\item[{1 }] symbol at particle's birth point
(at primary source)
\item[{2 }] symbol at the point of an electron
impact collision event
\item[{3 }] symbol at the point of a hard
elastic collision event
\item[{4 }] symbol at the point of a
charge exchange event
\item[{5 }] symbol at the point of a soft
elastic collision event
\item[{6 }] symbol at an intersection with
a ``non-default" or ``additional" surface
\item[{7 }] symbol at a
non-analog particle splitting point
\item[{8 }] symbol at a
non-analog particle killing point

(``Russian Roulette")
\item[{9 }]
symbol at a periodicity surface intersection point
\item[{10 }]
symbol at restart after splitting
\item[{11 }]
symbol at collision point saved for conditional exp. estimator
\item[{12 }]
symbol at a continuation of track for conditional exp. estimator
\item[{13 }]
symbol at a particle stopped at time limit (t-dep. mode)
\item[{14 }]
symbol at a particle stopped at generation limit (t-dep. mode)
\item[{15 }]
If an error is detected, by default a symbol is always plotted.
Plotting this symbol cannot be abandoned.
\end{list}
\item[\textbf{ILINIE }] ~
\begin{list}{ }{}
\item[{$\ne 0$ }]
connect two successive events by a straight line.
If a continuation of a track is computed for the conditional
expectation estimators, this part is represented by a dotted line,
see eq. (1.2.5).
\item[{= 0 }]
only symbols (if any selected) will be plotted along the history
\end{list}
\end{list}

\textsl{* 11B3 Graphical output of volume averaged tallies}

\begin{list}{ }{}
\item[\textbf{NVOLPL }]
Total number of pictures from volume averaged tallies.
\item[\textbf{PLTSRC(ISTRA), (ISTRA=0,NSTRAI) }] ~

Profiles of the selected
volumetric tallies are plotted
for those strata for which PLTSRC(ISTRA) = TRUE. In case
PLTSRC(0) = TRUE, the summation over all strata is plotted for
these tallies.

Note: graphical output for individual strata is only possible if the
data for strata have been saved on file (NFILE-N=1,2 option, input block 1).
Otherwise only the last stratum currently on storage arrays (mostly: sum over strata) is available.
\item[\textbf{NSPTAL }]
Number of different ``species" (quantities) to be plotted for this
tally into one picture.

For 3D plots, NSPTAL = 1, except for the
LVECT3 or LRPVC3 vector field options, in which case NSPTAL = 2. See below.
\item[\textbf{PLTL2D }]
Switch to activate a 2D plot
\item[\textbf{PLTL3D }]
Switch to activate a 3D plot
\item[\textbf{PLTLLG }]
The tally is plotted on a logarithmic scale.
\item[\textbf{PLTLER }]
Standard deviation is plotted if available. In case of
2D plot error bars are plotted along the curves. In case of 3D plot
a full standard deviation profile is plotted on a separate picture.
\item[\textbf{TALZMI }]
Minimum ordinate value for the plot.

If TALZMI = 666.0 then the minimum is searched for in the data array.

Even if PLTLLG (log. scale), the true
minimum ordinate value (not the logarithm thereof) must be specified.

If TALZMI = 666.0 and PLTLLG, ZMIN=1.E-48 is used as cut-off value.
\item[\textbf{TALZMA }]
like TALZMI, but
maximum ordinate value for the plot.

If TALZMA = 666.0 then the maximum is searched for in the data array.

Even if PLTLLG (log. scale), the true
maximum ordinate value (not the logarithm thereof) must be specified.
\end{list}

The following input variables are needed only if PLTL2D = TRUE

\begin{list}{ }{}
\item[\textbf{TALXMI = TALXMA = 0.0 }]
the radial (x-) grid is used as abscissa
for plotting (default).
\end{list}
The next options allow to change the x axis, e.g. for cases when a meaningful
radial (x-) grid does not exist for the curve to be plotted. No interpolation to a ``true" radial (x-) grid
is involved, so the resulting curves may change shape on the plot depending on options chosen here.
\begin{list}{ }{}
\item[\textbf{TALXMI $\neq$ TALXMA }] ~

if TALXMI $<$ TALXMA, then a grid equidistant on a linear scale is used,
TALXMI is then the minimum abscissa value for the plot,
TALXMA is then the maximum abscissa value for the plot.

if TALXMI $>$ TALXMA and both are positive, then a grid equidistant on a
log. scale is used.
TALXMI is then the minimum abscissa value for the plot,
TALXMA is then the maximum abscissa value for the plot.

if TALXMI $>$ TALXMA and at least one of them is negative,
then a user defined grid XXP2D\_USR is used. This must have been
defined in one of the user-routines for this figure no. N=IBLD,
i.e., the array XXP2D(J,IBLD) must have been set, J the grid index.
\item[\textbf{ISPZTL }]
Index of the species of the tally that will be plotted.

If ISPZTL = 0, the tally obtained by summation over its
species index is plotted.
\item[\textbf{NPTALI }]
Index of tally to be plotted.
\item[\textbf{NPLIN2 }]
Index of the first cell that is displayed on the plot,
for 2D plot only.
\item[\textbf{NPLOT2 }]
Index of the last cell that is displayed on the plot,
for 2D plot only.
\item[\textbf{NPLDL2 }]
Increment for the cell indices,
for 2D plot only.
\end{list}

The following input variables are needed only if PLTL3D = TRUE

\begin{list}{ }{}
\item[\textbf{LHIST3 }]
make a 3D histogram plot.
\item[\textbf{LCNTR3 }]
make a contour plot.
\item[\textbf{LSMOT3 }]
make a surface plot in a 3D cube.
\item[\textbf{LRAPS3 }]
produce files for RAPS graphics system. If additional surfaces are
to be indicated in these blocks, PLSTOR =.TRUE., see sub-block 11B1.
\item[\textbf{LVECT3}]
make a vector field plot.

NSPTAL must be 2 in this case. The first of these
2 tallies is taken as field of x - coordinates of the vector field to
to be plotted, and the second tally is the y - coordinate field.
\item[\textbf{LRPVC3}]
produce files for RAPS graphics system for vector field plot.
Co-ordinates of vectors as in LVECT3 option. If additional
surfaces are to be indicated in these blocks, PLSTOR=TRUE, see
sub-block 11B1.

\item[\textbf{LPRAD3 }]
radial (x) co-ordinate is fixed. Plot tally versus poloidal (y) and toroidal (z)
co-ordinate. See IPROJ3 flag below.
\item[\textbf{LPPOL3 }]
poloidal (y) co-ordinate is fixed. Plot tally versus radial (x) and toroidal (z)
co-ordinate
\item[\textbf{LPTOR3 }]
toroidal (z) co-ordinate is fixed. Plot tally versus radial (x) and poloidal (y)
co-ordinate
\item[\textbf{ISPTAL }]
species index of tally to be plotted. ISPTAL = 0 means: sum over species
index.
\item[\textbf{NPTALI }]
Index of tally to be plotted.
\item[\textbf{IPROJ3 }]
in case of 3d calculation, IPROJ3 is the index of
the mesh cell of the grid, for which the corresponding
co-ordinate is fixed. (See LPRAD3, LPPLO3, LPTOR3 flags above).
\item[\textbf{NPLI13 }]
\item[\textbf{NPL013 }]
\item[\textbf{NPLI23 }]
\item[\textbf{NPLO23 }]
\item[\textbf{TALW1 }]
first viewing angle for 3D plot
\item[\textbf{TALW2 }]
second viewing angle for 3D plot
\item[\textbf{FCABS1 }]
\item[\textbf{FCABS2 }]
\end{list}
